\documentclass[../../Bayesian_Inference.tex]{subfiles}

\begin{document}

\section{General Framework of Bayesian Regression}

Consider a training dataset
\[
    \mathcal{D} = \left\{(\bm{x}_i,y_i)\right\}_{i=1}^{N},
\]
consisting of $N$ observations, where $\bm{x}_i \in \mathbb{R}^{p}$ denotes the $p$-dimensional predictor vector associated with the $i$-th observation, and $y_i \in \mathbb{R}$ denotes the corresponding scalar response. Define the design matrix $\mathbf{X}$, response vector $\bm{y}$, and regression coefficient vector $\bm{w}$ as
\[
    \mathbf{X} =
    \begin{bmatrix}
        \bm{x}_1^{\top} \\
        \vdots \\
        \bm{x}_N^{\top}
    \end{bmatrix}
    \in \mathbb{R}^{N \times p},
    \qquad
    \bm{y} =
    \begin{bmatrix}
        y_1 \\
        \vdots \\
        y_N
    \end{bmatrix}
    \in \mathbb{R}^{N},
    \qquad
    \bm{w} =
    \begin{bmatrix}
        w_1 \\
        \vdots \\
        w_p
    \end{bmatrix}
    \in \mathbb{R}^{p}.
\]

Given the predictor vector $\bm{x}_i$ and the regression coefficient vector $\bm{w}$, the response $y_i$ is characterized by the conditional distribution
\[
    p(y_i \mid \bm{x}_i, \bm{w}).
\]
Assuming that the observations are conditionally independent given $\bm{w}$, the likelihood of the observed response vector is
\[
    p(\bm{y} \mid \mathbf{X}, \bm{w})
    =
    \prod_{i=1}^{N}
    p(y_i \mid \bm{x}_i, \bm{w}).    
\]

In the Bayesian framework, the regression coefficient vector $\bm{w}$ is treated as a random variable and assigned a prior distribution $p(\bm{w})$. Given the observed dataset $\mathcal{D} = (\mathbf{X},\bm{y})$, Bayes' theorem gives
\[
    p(\bm{w} \mid \mathcal{D})
    =
    p(\bm{w} \mid \mathbf{X}, \bm{y})
    =
    \frac{p(\bm{w},\mathbf{X},\bm{y})}
         {p(\mathbf{X},\bm{y})}.
\]
Using
\[
    p(\mathbf{X},\bm{y})
    =
    p(\bm{y}\mid\mathbf{X})p(\mathbf{X}),
\]
the posterior distribution can be written as
\[
    p(\bm{w} \mid \mathbf{X}, \bm{y})
    =
    \frac{p(\bm{w},\bm{y}\mid\mathbf{X})}
         {p(\bm{y}\mid\mathbf{X})}
    =
    \frac{
        p(\bm{y}\mid\mathbf{X},\bm{w})
        p(\bm{w}\mid\mathbf{X})
    }{
        p(\bm{y}\mid\mathbf{X})
    }.
\]
Assuming that the prior distribution of $\bm{w}$ is independent of the design matrix $\mathbf{X}$, i.e.,
\[
    p(\bm{w}\mid\mathbf{X}) = p(\bm{w}),
\]
the posterior distribution becomes
\begin{equation}
    \boxed{p(\bm{w} \mid \mathbf{X}, \bm{y})
    =
    \frac{
        p(\bm{y} \mid \mathbf{X}, \bm{w})
        p(\bm{w})
    }{
        p(\bm{y}\mid\mathbf{X})
    }}.    
\end{equation} 
The denominator
\begin{equation}
    \boxed{p(\bm{y}\mid\mathbf{X})
    =
    \int
    p(\bm{y}\mid\mathbf{X},\bm{w})
    p(\bm{w})
    \mathrm{d}\bm{w}}    
\end{equation} 
is referred to as the marginal likelihood, or model evidence. 

Since it does not depend on $\bm{w}$, the posterior distribution can be expressed up to a proportionality constant as
\begin{equation}
    \boxed{p(\bm{w} \mid \mathbf{X}, \bm{y})
    \propto
    p(\bm{y} \mid \mathbf{X}, \bm{w})
    p(\bm{w})},   
\end{equation}
where $p(\bm{y} \mid \mathbf{X}, \bm{w})$ represents the likelihood and $p(\bm{w})$ represents the prior distribution. Therefore, the fundamental principle of Bayesian inference can be summarized as
\[
    \boxed{\text{Posterior}
    \propto
    \text{Likelihood}
    \times
    \text{Prior}}.
\]

The objective of Bayesian regression is to infer the posterior distribution of the regression coefficient vector $\bm{w}$ given the observed dataset $\mathcal{D}$, and subsequently to obtain the posterior predictive distribution of the response $y^{*}$ for a new observation with predictor vector $\bm{x}^{*}$.
\[
    p(\bm{y}^* | \mathbf{X}^{*}, \mathcal{D}) = \int p(\bm{y}^* | \mathbf{X}^{*}, \bm{w}) p(\bm{w} | \mathcal{D}) \mathrm{d} \bm{w}.
\]



\end{document}